\chapter{\IfLanguageName{dutch}{Evaluatie en validatie van het prototype}{Evaluation and validation of the prototype}}%
\label{ch:evaluatie-prototype}

\section{Inleiding}
\label{sec:eval-inleiding}

Dit hoofdstuk evalueert het prototype dat is ontwikkeld op basis van de ontwerprichtlijnen uit het vorige hoofdstuk. De evaluatie focust op drie dimensies: functionele correctheid, conformiteit met de geformuleerde ontwerprichtlijnen, en bruikbaarheid van de gebruikersinterface.

\subsection{Afbakening en motivering van de scope}

De tool die in dit hoofdstuk wordt geëvalueerd, ondersteunt op dit moment één back-up-methode: incrementele tar+zstd. Btrfs-snapshotting werd in het finale prototype weggelaten. De doorslaggevende reden is: het aandeel computergebruikers dat een btrfs-bestandssysteem gebruikt is bijzonder laag, waardoor de relevantie van deze methode voor de doelgroep (KSP-spelers in het algemeen) beperkt blijft. Btrfs-ondersteuning wordt om die reden beschouwd als een uitbreiding voor toekomstig werk, detectie voor het type bestandssysteem is wel al geïmplementeerd. Sectie~\ref{sec:eval-conformiteit} gaat hier dieper op in bij de bespreking van de eerste ontwerprichtlijn.

Een tweede afbakening betreft de evaluatiemethodologie. De UX-evaluatie in sectie~\ref{sec:eval-heuristiek} steunt op een heuristische analyse volgens de tien usability-principes van Nielsen, uitgevoerd door de ontwikkelaar zelf op basis van praktisch gebruik van het prototype op een werkelijke KSP-installatie met meer dan 100 mods en verschillende savegames.

\subsection{Keuze van de programmeertaal}
\label{subsec:eval-taal}

Voor de implementatie van het prototype werd Python gekozen. Deze keuze verdient toelichting, omdat CKAN (de bestaande modlauncher voor KSP) in C\# is geschreven. Vier overwegingen lagen aan de basis van de keuze voor Python in deze fase:

\begin{itemize}
    \item \textbf{Aard van de werklast.} De kernoperaties van de tool (\texttt{btrfs subvolume snapshot}, \texttt{tar}, \texttt{zstd}, \texttt{rsync}, \texttt{findmnt}) zijn allemaal externe binaries die via shell-commando's worden aangeroepen. De programmeertaal speelt hier geen rol in de prestaties, alle rekenintensieve werk gebeurt in de aangeroepen tools. Pythons \texttt{subprocess}-module biedt de meest directe en leesbare integratie voor dit soort patronen op Linux.
    \item \textbf{Snelheid van ontwikkeling.} Voor een proof of concept weegt ontwikkelsnelheid zwaarder dan distributiegemak. Pythons rijke standaardbibliotheek (\texttt{json}, \texttt{pathlib}, \texttt{dataclasses}, \texttt{subprocess}) en PySide6 voor Qt laten toe om in een beperkte tijd tot een werkende GUI te komen.
    \item \textbf{Distributievereisten zijn beperkt voor een PoC.} Het belangrijkste voordeel dat C\# zou bieden, een single- binary distributie via .NET, is voor een prototype dat op één testsysteem moet draaien geen doorslaggevend criterium.
    \item \textbf{CKAN-integratie als toekomstig werk.} Een rechtstreekse integratie van het back-upsysteem in CKAN is het ultieme einddoel van dit type tool, maar valt buiten de scope van dit prototype. Wanneer die integratie in de toekomst wordt ondernomen, is een C\#-implementatie de natuurlijke keuze. Voor de huidige fase, waarin de focus ligt op het valideren van de ontwerprichtlijnen via een werkende standalone tool, biedt Python de beste route.
\end{itemize}

De combinatie van Python met PySide6 (Qt for Python) levert een moderne dark-themed GUI op die past bij hedendaagse desktops. Figuur~\ref{fig:eval-home} toont het hoofdscherm van het prototype waarin een KSP-installatie met 113 mods is geconfigureerd.

\begin{figure}[h]
    \centering
    \includegraphics[width=0.95\linewidth]{home_screen.png}
    \caption{Configuratie-scherm van het prototype, met een werkelijke KSP mod-installatie van 20{,}5~GB en 113 mods.}
    \label{fig:eval-home}
\end{figure}


\section{Functionele evaluatie via testmatrix}
\label{sec:eval-functioneel}

Om de functionele correctheid van het prototype systematisch te beoordelen, werd een testmatrix opgesteld die elk afzonderlijk aspect van de tool valideert. De testen zijn gegroepeerd in vier clusters: basis back-up en restore, selectie en interface, configuratie en metadata, en edge cases met foutafhandeling. Elk testgeval beschrijft een doel, een uitgevoerde stap en een verwacht resultaat, waarna het werkelijke gedrag wordt gerapporteerd.

\subsection{Cluster 1: Basis back-up en restore}

Tabel~\ref{tab:eval-cluster1} toont de testen van de kernfunctionaliteiten. Deze testen werden uitgevoerd op een werkelijke KSP-installatie met 113 mods die 20{,}5~GB omvatten. Eén van de testen (volledige back-up van de mod-set) leverde meteen ook een nuttige indicatieve meting op: de back-up van alle 113 mods (level~0) duurde ongeveer 37{,}5~seconden en resulteerde in een gecomprimeerd archief van 10{,}3~GB. Een volledige restore van die set duurde ongeveer 25~ seconden. Deze cijfers fungeren in dit hoofdstuk niet als benchmark, maar als bewijs dat het prototype daadwerkelijk werkt op realistische werklasten zonder de UI te bevriezen.

\begin{table}[h]
    \centering
    \caption{Cluster 1: basis back-up en restore.}
    \label{tab:eval-cluster1}
    \small
    \begin{tabular}{p{0.32\linewidth}p{0.42\linewidth}p{0.16\linewidth}}
        \toprule
        \textbf{Testgeval} & \textbf{Verwacht resultaat} & \textbf{Status} \\
        \midrule
        Eerste back-up van een save als level 0 & Volledig basisarchief wordt aangemaakt en metadata-entry met \texttt{increment\_level=0}. & Geslaagd \\
        Eerste back-up van een mod als level 0 & Identiek aan vorige test, maar voor een mod-map onder \texttt{GameData/}. & Geslaagd \\
        Tweede back-up van hetzelfde item creëert level 1 & Tar gebruikt het bestaande snar-bestand en produceert enkel de wijzigingen sinds de vorige back-up. & Geslaagd \\
        Restore van een level-0-punt & Werkmap wordt teruggebracht naar de toestand van het basisarchief. & Geslaagd \\
        Restore van een level-2-punt & Drie archieven (level 0, 1 en 2) worden sequentieel uitgepakt; de werkmap reflecteert de cumulatieve toestand. & Geslaagd \\
        Multi-item back-up (saves en mods) in één actie & Alle geselecteerde items krijgen dezelfde \texttt{set\_id} en worden gegroepeerd weergegeven in de restore-lijst. & Geslaagd \\
        Back-up en restore van een grote werklast (20{,}5~GB, 113 mods) & Operatie voltooit zonder UI-bevriezing; gemeten $\sim$37{,}5~s back-up, $\sim$25~s restore. & Geslaagd \\
        \bottomrule
    \end{tabular}
\end{table}

Figuur~\ref{fig:eval-backup-progress} toont het back-up-scherm tijdens een volledige back-up van de 113 mods. De vooruitgangsbalk en de live log onderaan houden de gebruiker geïnformeerd zonder dat de UI bevriest. De aangeroepen \texttt{tar}-processen draaien op een achtergrondthread via Qt's \texttt{QThread}-mechanisme.

\begin{figure}[h]
    \centering
    \includegraphics[width=0.95\linewidth]{backup_scherm_tijdens_backup.png}
    \caption{Back-up-scherm tijdens uitvoering. De progressbar toont dat de operaties nog bezig zijn en de log onderaan rapporteert per item de status en grootte. De gebruikersinterface blijft responsief tijdens de hele operatie.}
    \label{fig:eval-backup-progress}
\end{figure}

\subsection{Cluster 2: Selectie en interface}

Tabel~\ref{tab:eval-cluster2} verzamelt de testen rond gebruikersinteractie en visuele weergave.

\begin{table}[h]
    \centering
    \caption{Cluster 2: selectie en interface.}
    \label{tab:eval-cluster2}
    \small
    \begin{tabular}{p{0.32\linewidth}p{0.42\linewidth}p{0.16\linewidth}}
        \toprule
        \textbf{Testgeval} & \textbf{Verwacht resultaat} & \textbf{Status} \\
        \midrule
        Selectie van meerdere saves uit een lijst & Enkel de aangevinkte saves worden in de back-up opgenomen. & Geslaagd \\
        Stock-content (\texttt{Squad}, \texttt{SquadExpansion}) wordt overgeslagen & Deze mappen verschijnen niet in de mod-lijst. & Geslaagd \\
        Multi-select restore over verschillende back-up-sets & Items uit verschillende sets kunnen tegelijk geselecteerd en in één actie hersteld worden. & Geslaagd \\
        Boomstructuur groepeert per back-up-actie en daarbinnen per type & Restore-tab toont sets met daaronder ``Saves''- en ``Mods''-groepen. & Geslaagd \\
        Groene cirkel toont het huidige back-uppunt & Na een back-up of restore krijgt het corresponderende back-uppunt een groene markering; oudere punten van hetzelfde item verliezen die markering. & Geslaagd \\
        \bottomrule
    \end{tabular}
\end{table}

Figuur~\ref{fig:eval-backup-screen} illustreert het selectiescherm na voltooiing van een back-up van alle 113 mods. De saves- en mods-kolommen tonen aanvinkbare lijsten, stock-mappen zoals \texttt{Squad} ontbreken. Figuur~\ref{fig:eval-restore-after} toont de restore-tab na een herstel van een save naar level 2: de groene cirkel naast het zonet herstelde punt geeft duidelijk aan welk punt op dit moment actief is in de werkmap.

\begin{figure}[h]
    \centering
    \includegraphics[width=0.95\linewidth]{backup_screen.png}
    \caption{Back-up-scherm met selectie-interface. Saves en mods staan in afzonderlijke kolommen. De methode (incrementele tar+zstd) wordt rechtsbovenin altijd getoond.}
    \label{fig:eval-backup-screen}
\end{figure}

\begin{figure}[h]
    \centering
    \includegraphics[width=0.95\linewidth]{restore_screen_na_restore_van_level_3.png}
    \caption{Restore-tab na een herstel naar een level-2-punt. De log onderaan toont dat drie archieven sequentieel werden uitgepakt. De groene cirkel naast het herstelde punt (level 2) markeert het als de huidige toestand; de oudere punten van datzelfde item zijn niet meer gemarkeerd.}
    \label{fig:eval-restore-after}
\end{figure}

\subsection{Cluster 3: Configuratie en metadata}

Tabel~\ref{tab:eval-cluster3} verzamelt de testen rond configuratie en metadata. De testen in dit cluster tonen aan dat de scheiding tussen globale en per-installatie state correct werkt. De globale configuratie in \texttt{\textasciitilde/.config/ksp-backup/installations.json} beperkt zich bewust tot een lijst van paden naar bekende installaties. Alle eigenlijke back-upmetadata staat in de \texttt{metadata.json} binnen elke KSP-installatie zelf.


\begin{table}[h]
    \centering
    \caption{Cluster 3: configuratie en metadata.}
    \label{tab:eval-cluster3}
    \small
    \begin{tabular}{p{0.32\linewidth}p{0.42\linewidth}p{0.16\linewidth}}
        \toprule
        \textbf{Testgeval} & \textbf{Verwacht resultaat} & \textbf{Status} \\
        \midrule
        Toevoegen van een KSP-installatie persisteert na herstart & De installatie verschijnt opnieuw in de lijst na het sluiten en opnieuw openen van de applicatie. & Geslaagd \\
        Verwijderen van een installatie raakt de bestanden niet aan & De installatie wordt enkel uit de globale lijst gehaald; de KSP-folder en \texttt{Backups/}-submap blijven intact. & Geslaagd \\
        Beheren van meerdere installaties tegelijk & Verschillende KSP-installaties kunnen onafhankelijk worden geconfigureerd; elke installatie heeft zijn eigen \texttt{metadata.json}. & Geslaagd \\
        Detectie van het bestandssysteem & \texttt{findmnt} geeft \texttt{btrfs} terug voor de getoonde installatie; de configuratie-tab rapporteert dit correct (zie figuur~\ref{fig:eval-home}). & Geslaagd \\
        \bottomrule
    \end{tabular}
\end{table}

\subsection{Cluster 4: Edge cases en foutafhandeling}

\begin{table}[h]
    \centering
    \caption{Cluster 4: edge cases en foutafhandeling.}
    \label{tab:eval-cluster4}
    \small
    \begin{tabular}{p{0.32\linewidth}p{0.42\linewidth}p{0.16\linewidth}}
        \toprule
        \textbf{Testgeval} & \textbf{Verwacht resultaat} & \textbf{Status} \\
        \midrule
        Verwijderen van een level 0 leidt tot cascade-cleanup & Het basisarchief, alle latere incrementen van datzelfde item én het snar-bestand worden samen verwijderd. & Geslaagd \\
        Verwijderen van enkel een tussenliggende increment behoudt de chain & Het basisarchief en eventueel andere incrementen blijven bruikbaar; het snar-bestand blijft eveneens behouden. & Geslaagd \\
        Bevestigingsdialoog vóór restore en delete & Een dialoog toont alle te overschrijven of te verwijderen items vóór de uitvoering en vraagt expliciet om bevestiging. & Geslaagd \\
        Cascade-waarschuwing bij verwijderen van een level 0 & De delete-dialoog vermeldt expliciet welke afhankelijke incrementen mee zullen worden verwijderd. & Geslaagd \\
        \bottomrule
    \end{tabular}
\end{table}

Tabel~\ref{tab:eval-cluster4} toont de testen van de edge cases en foutafhandeling. Bij het opzetten van deze testmatrix kwam tijdens de ontwikkeling een bug aan het licht: bij het verwijderen van een level-0-basisarchief werd het snar-bestand niet mee opgeruimd, waardoor de eerstvolgende back-up van datzelfde item als level 1 werd aangemaakt op basis van een verouderd snar-bestand. De fout werd verholpen door bij elke verwijdering te controleren of er nog een level-0-archief overblijft voor het betreffende item, en zo niet, het snar-bestand én alle verweesde incrementen mee op te ruimen. Figuur~\ref{fig:eval-delete-dialog} toont de bevestigingsdialoog die de gebruiker informeert over dit scenario.

\begin{figure}[h]
    \centering
    \includegraphics[width=0.55\linewidth]{verwijderen_van_een_level_0.png}
    \caption{Bevestigingsdialoog bij het verwijderen van een level-0-basisarchief. De dialoog toont expliciet welke afhankelijke incrementen mee zullen worden opgeruimd, zodat de gebruiker geen verrassingen krijgt na bevestiging.}
    \label{fig:eval-delete-dialog}
\end{figure}

\section{Conformiteitscheck tegen de ontwerprichtlijnen}
\label{sec:eval-conformiteit}

In het vorige hoofdstuk werden elf ontwerprichtlijnen geformuleerd op basis van de technische analyse. Tabel~\ref{tab:eval-conformiteit} geeft een beknopt overzicht van de implementatiestatus van elke richtlijn in het huidige prototype. De daaropvolgende toelichting behandelt enkel de richtlijnen die nuance vereisen.

\begin{table}[h]
    \centering
    \caption{Conformiteit van het prototype met de ontwerprichtlijnen.}
    \label{tab:eval-conformiteit}
    \small
    \begin{tabular}{p{0.04\linewidth}p{0.40\linewidth}p{0.18\linewidth}p{0.25\linewidth}}
        \toprule
        \textbf{\#} & \textbf{Richtlijn} & \textbf{Status} & \textbf{Toelichting} \\
        \midrule
        1 & Snapshot primair op btrfs              & Niet geïmplementeerd & Uitbreiding \\
        2 & Incrementele tar+zstd als alternatief  & Volledig             & Kernimplementatie \\
        3 & Differentieel afgeraden                & Conform              & Geen optie aangeboden \\
        4 & Volledige back-up zelden zinvol        & Conform              & Level 0 dient enkel als basis \\
        5 & \texttt{zstd} standaard                & Volledig             & Kernimplementatie \\
        6 & \texttt{xz} optioneel voor max. opslag & Niet geïmplementeerd & Uitbreiding \\
        7 & Andere algoritmes afgeraden            & Conform              & Geen optie aangeboden \\
        9 & Lokaal werken, geen netwerk            & Volledig             & Geen netwerklogica \\
        10 & Methode en compressie orthogonaal     & Conform              & Architectuur scheidt beide \\
        11 & Voorspelbaarheid als ontwerpkwaliteit & Volledig             & Methode zichtbaar in UI, sets gegroepeerd, huidige status zichtbaar \\
        \bottomrule
    \end{tabular}
\end{table}

\paragraph{Richtlijn 1: Snapshotting als uitbreiding.}
De technische analyse heeft btrfs-snapshotting aangeduid als de primaire methode voor opstellingen die het ondersteunen. In de finale tool werd deze methode echter weggelaten. Het percentage computergebruikers dat een btrfs-bestandssysteem gebruikt is bijzonder laag, en de typische KSP-speler bevindt zich op Windows met NTFS of op een Linux-systeem met ext4. Snapshotting heeft daardoor een beperkte praktische relevantie voor de doelgroep, ondanks de gunstige technische eigenschappen. Btrfs-snapshotting wordt om die reden behandeld als een uitbreiding in plaats van een kernvereiste. Het prototype valt voor alle bestandssystemen terug op de tweede aanbeveling uit de analyse: incrementele tar+zstd.

\paragraph{Richtlijn 6: \texttt{xz} als uitbreiding.}
De tool gebruikt momenteel uitsluitend \texttt{zstd} als compressie-algoritme, conform richtlijn~5. Een optionele variant met \texttt{xz} voor gebruikers die maximale opslagbesparing wensen werd nog niet geïmplementeerd. Het toevoegen van een tweede algoritme-keuze in de configuratie-tab is technisch eenvoudig en wordt beschouwd als een gepaste uitbreiding.


\section{Heuristische evaluatie van de gebruikersinterface}
\label{sec:eval-heuristiek}

De bruikbaarheid van de UI werd beoordeeld via een heuristische evaluatie volgens de tien principes van Nielsen. Per principe wordt kort vermeld hoe het prototype ermee omgaat, met een lichte status (sterk, voldoende of aandachtspunt).

\begin{description}
    \item[Zichtbaarheid van systeemstatus: sterk.] De progressbar en live log tijdens een back-up of restore tonen duidelijk welke operatie aan de gang is. De groene cirkel naast het huidige back-uppunt informeert over de actieve toestand. De gekozen back-upmethode wordt rechtsbovenin de back-up-tab permanent zichtbaar gehouden.
    \item[Match met de echte wereld: sterk.] Begrippen als ``back-up maken'', ``herstellen'' en ``installatie toevoegen'' sluiten aan bij gangbaar taalgebruik van eindgebruikers. Technische termen zoals \texttt{btrfs} of \texttt{snar} blijven achter de schermen.
    \item[Gebruikerscontrole en vrijheid: voldoende.] Bevestigingsdialogen vóór restore en delete voorkomen accidentele onomkeerbare acties. Een ``ongedaan maken''-functie ontbreekt: eens een restore is uitgevoerd, kan de gebruiker enkel terugvallen op een ander back-uppunt.
    \item[Consistentie en standaarden: sterk.] De Qt/KDE- look-and-feel zorgt voor herkenbaar gedrag van knoppen, dialogen en lijsten. Acties die data wijzigen krijgen consistent een blauwe (primaire) of rode (gevaarlijke) accentkleur.
    \item[Foutpreventie: sterk.] Stock-content wordt automatisch uitgesloten van de mod-lijst, waardoor de gebruiker geen onnodige back-ups kan maken van data die niet vervangen mag worden. De cascade-waarschuwing bij het verwijderen van een level 0 (figuur~\ref{fig:eval-delete-dialog}) toont expliciet de afhankelijkheden vóór bevestiging.
    \item[Herkennen in plaats van onthouden: voldoende.] Back-uppunten zijn gegroepeerd per back-up-actie en daarbinnen per type. De gebruiker hoeft geen IDs of timestamps te onthouden, het visueel herkennen van een set en het aanvinken van items volstaat. De mogelijkheid om een naam te geven aan een back-up ontbreekt nog.
    \item[Flexibiliteit en efficiëntie: voldoende.] Knoppen voor ``alles selecteren'' en ``niets selecteren'' versnellen bulk-operaties. De boomstructuur kan in zijn geheel uitgevouwen of opgevouwen worden. Sneltoetsen voor frequente acties ontbreken.
    \item[Esthetisch en minimalistisch ontwerp — sterk.] De UI bestaat uit drie tabs zonder overbodige elementen. Hint-teksten verschijnen enkel waar de gebruiker context nodig heeft.
    \item[Foutdiagnose: voldoende.] Foutmeldingen verschijnen in de log onderaan en in een dialoog. Voor een PoC is dit voldoende, maar uitgebreidere foutclassificatie (bijvoorbeeld onderscheid tussen ``schijf vol'' en ``rechten ontbreken'') zou de gebruiksvriendelijkheid in productie verhogen.
    \item[Hulp en documentatie: aandachtspunt.] De tool bevat enkel hint-labels per tab. Een aparte help-pagina, tooltip-uitleg op individuele velden of een korte introductie-tour ontbreekt. Voor een PoC is de tool eenvoudig genoeg om dit te aanvaarden, maar voor productiekwaliteit is documentatie een nodige toevoeging.
\end{description}

Een heuristische evaluatie door de ontwikkelaar zelf vangt blinde vlekken niet op die echte gebruikers wel zouden ontdekken. Deze sectie biedt daarom een gestructureerde reflectie op de UI-kwaliteit, geen gevalideerde usability-meting. Voor een productieversie zou gebruikersonderzoek met effectieve KSP-spelers een logische volgende stap zijn.


\section{Beperkingen en aanknopingspunten voor toekomstig werk}
\label{sec:eval-toekomst}

Onderstaande punten geven aan welke aspecten buiten scope zijn gehouden en hoe zij in toekomstig werk kunnen worden aangepakt.

\paragraph{Geen evaluatie met externe testers.}
De UX-evaluatie steunt op een heuristische analyse door de ontwikkelaar zelf. Een gevalideerde usability-studie met effectieve KSP-spelers zou pas mogelijk zijn na een eerste distributie van de tool en valt buiten de scope van deze bachelorproef.

\paragraph{Btrfs-snapshotting.}
Zoals besproken in sectie~\ref{sec:eval-conformiteit} werd de snapshot-methode niet geïmplementeerd in de finale tool wegens de beperkte praktische relevantie voor de doelgroep. Voor gebruikers op een btrfs-systeem die het maximale uit hun bestandssysteem willen halen, blijft dit een goede uitbreiding.

\paragraph{\texttt{xz} als optionele compressie.}
De tool gebruikt momenteel uitsluitend \texttt{zstd}. Een configuratie-optie waarmee de gebruiker kan kiezen voor \texttt{xz} wanneer maximale opslagbesparing is gewild zou conform richtlijn~6 zijn en is technisch eenvoudig toe te voegen.

\paragraph{automatische opruim.}
Het prototype voorziet geen automatisch beheer van oude back-uppunten. Bij langdurig gebruik kan de hoeveelheid back-uppunten oplopen. Een mechanisme dat oudere punten verwijdert volgens een instelbare  optie, zou de opslagvoetafdruk op lange termijn beheersbaar houden.

\paragraph{Scheduling van automatische back-ups.}
Op dit moment maakt de gebruiker manueel een back-up via de back-up-tab. Een planning-mechanisme dat automatisch een back-up maakt vlak vóór een mod-update of op vaste tijdstippen, sluit aan bij het oorspronkelijke gebruiksscenario en is een logische uitbreiding.

\paragraph{Multi-game-ondersteuning.}
Het prototype ondersteunt enkel KSP. De architectuur is generiek genoeg om uitgebreid te worden naar andere mod-gebaseerde games.

\paragraph{Integratie in een modlauncher.}
Het ultieme doel van een back-upsysteem zoals dit is rechtstreekse integratie in een bestaande modlauncher. Voor KSP betekent dit een integratie in CKAN. Aangezien CKAN geschreven is in C\#, vraagt deze integratie een herschrijving van de tool in c\#. Dit valt buiten de scope van het prototype en wordt beschouwd als de volgende stap naar een volwaardig back-up systeem.


\section{Samenvatting}
\label{sec:eval-samenvatting}

De evaluatie heeft het prototype op drie dimensies onderzocht: functionele correctheid, conformiteit met de ontwerprichtlijnen, en bruikbaarheid van de gebruikersinterface. De testmatrix bevestigt dat de tool zijn gedragsspecificatie nakomt over twintig testgevallen, gegroepeerd in vier clusters, inclusief edge cases zoals de cascade-cleanup bij het verwijderen van een basisarchief. De conformiteitscheck toont dat zes van de tien besproken richtlijnen volledig zijn geïmplementeerd, drie conform zijn afgehandeld door bewust geen optie aan te bieden, en twee (snapshotting en \texttt{xz}) als gemotiveerde uitbreiding voor toekomstig werk zijn aangemerkt. De heuristische UI-evaluatie identificeert het prototype als bruikbaar binnen de PoC-context, met enkele aandachtspunten waarvan de belangrijkste (documentatie, sneltoetsen, ``ongedaan maken'') typisch onder productiekwaliteit horen.

Een positieve observatie tijdens de evaluatie was dat het prototype probleemloos werkt op realistische werklasten: indicatieve metingen op een KSP-installatie van 20{,}5~GB met 113 mods leverden een back-uptijd van ongeveer 37{,}5~seconden en een restoretijd van ongeveer 25~seconden op, terwijl de UI gedurende de hele operatie responsief bleef. Het prototype demonstreert daarmee dat de in deze bachelorproef geformuleerde ontwerprichtlijnen niet alleen theoretisch verdedigbaar zijn, maar ook praktisch realiseerbaar in een eenvoudige Linux-applicatie die vandaag al bruikbaar is voor de oorspronkelijke gebruikssituatie.